\section{Memory Wall：模型增长快于数据供给能力}

上一章给出资源下界，本章聚焦课堂反复强调的 memory wall。Gholami 等人在 \emph{AI and Memory Wall} 中讨论了模型规模、memory capacity、bandwidth 与通信增长不匹配的问题。对 LLM 来说，training 需要参数、optimizer state、gradient 与 activation；inference 还要管理 KV cache。模型“放得下”与模型“喂得饱”是两个不同条件。

\subsection{容量墙、带宽墙与通信墙}

本节进一步把“memory wall”拆成三个可测问题：模型和运行状态能否放下、单位时间能否被送到 compute、扩展到多卡后能否高效同步。读下面的图时不要把右侧统称为“memory 不够”，而要分别记录 capacity、bandwidth 和 communication 的证据，因为三类问题需要不同的硬件和软件修复。

\lecturefigure{03-memory-wall.png}{Memory wall 的本质是 compute throughput 与 data supply 的速率失配。}{本地字幕 08:20--10:30；Gholami et al., \emph{AI and Memory Wall}。}

\begin{knowledgebox}{读图：先区分三种 wall}
\textbf{Capacity wall} 问参数、KV cache 和 activation 是否能放入 memory；\textbf{bandwidth wall} 问单位时间能否把数据送到 compute；\textbf{communication wall} 问多设备间同步和重分布是否吞掉扩展收益。三者可能同时存在，解决其中一个会把瓶颈暴露到下一层。
\end{knowledgebox}

\begin{importantbox}{Arithmetic intensity 决定“算力还是带宽”}
定义 arithmetic intensity
\[
I=\frac{F}{B_m},
\]
即每搬运一个 byte 完成多少次运算。若 $I$ 很低，即使 accelerator 峰值吞吐很高，kernel 也更可能受 memory bandwidth 限制；fusion、tiling、quantization 和 cache reuse 的价值，都可以理解为提高有效 $I$ 或减少 $B_m$。
\end{importantbox}

\subsection{术语消化：SRAM、DRAM、HBM 与 storage}

这组词经常被字幕和市场材料混为“memory”。读者应先掌握它们在 hierarchy 中的角色，再理解 Groq、NPU、GPU 或新架构的差异。

\begin{knowledgebox}{背景概念：四类存储层}
\begin{tabularx}{\textwidth}{L{0.17\textwidth}L{0.24\textwidth}X}
\toprule
术语 & 主要特征 & 在 AI 系统中的角色与限制 \\
\midrule
SRAM & Static Random-Access Memory，片上高速缓存，低延迟但面积昂贵 & 适合小而热的数据和确定性访问；容量小，不等于通用模型存储。 \\
DRAM & Dynamic Random-Access Memory，需要刷新，容量较大 & 主存或 accelerator 外部 memory；延迟和带宽受接口、封装与并发访问影响。 \\
HBM & High Bandwidth Memory，高带宽堆叠 DRAM & 通过宽接口提供高吞吐；容量、封装供给、成本和功耗仍是约束。 \\
SSD & 持久存储，容量大但延迟远高于 memory & 保存 checkpoint、模型 artifact 与冷数据；通常需要 staging，不能直接替代 HBM。 \\
\bottomrule
\end{tabularx}
\end{knowledgebox}

\lecturefigure{04-memory-hierarchy.png}{不同 memory tier 在 latency、bandwidth、capacity 和 energy 上交换。}{本地字幕 09:00--11:10；概念重绘。}

\begin{knowledgebox}{读图：模型 serving 要按“温度”放数据}
最热的 kernel state 和小型 buffer 留在 on-chip SRAM；活跃 weights 与 KV cache 争夺 HBM；更大的 host memory 承担 staging；SSD 保存冷模型和 checkpoint。若 routing 频繁切换模型，冷启动和搬运会成为主要成本，因此 model portfolio 设计必须与 memory placement 联动。
\end{knowledgebox}

\begin{warningbox}{“某芯片只适合文本”不是稳定结论}
课堂用具体公司快速说明 SRAM/DRAM/HBM 取舍，但产品能力会随 compiler、model support、memory topology 和 serving software 演化。讲义保留 hierarchy 原理，不把 2025 年现场判断当作永久产品分类。
\end{warningbox}

\subsection{本章小结}

Memory wall 不是单一器件短缺，而是 capacity、bandwidth、communication 与 software locality 的组合问题。选择芯片之前，应先量化模型的 working set、bytes per token、KV cache 增长与跨设备 traffic。

